\documentclass[12pt]{article}

\input{../../_assets/preambulo.tex}


\begin{document}

    % 1. Foto de fondo
    % 2. Título
    % 3. Encabezado Izquierdo
    % 4. Color de fondo
    % 5. Coord x del titulo
    % 6. Coord y del titulo
    % 7. Fecha

    
    \input{../../_assets/portada}
    \portadaExamen{ffccA4.jpg}{Álgebra II\\Examen IX}{Álgebra II. Examen IX}{MidnightBlue}{-8}{28}{2025-26}{David Muñoz Gómez}

    \begin{description}
        \item[Asignatura] Álgebra II.
        \item[Curso Académico] 25/26.
        \item[Grado] DGIIM.
        \item[Grupo] Único.
        \item[Profesor] Aurora Inés del Río Cabeza.
        \item[Descripción] Prueba intermedia II.
        \item[Fecha] 1/06/2026.
        % \item[Duración] ---.
    \end{description}
    \newpage


    % ------------------------------------
    
    \begin{ejercicio}[1.0 punto]
        \textbf{Cuestionario}

        En las siguientes cuestiones responde ``VERDADERO'' o ``FALSO'' y haz un breve razonamiento para justificar tu respuesta. Cada respuesta correcta puntúa 0,1 puntos.

        \begin{enumerate}[label=\arabic*)]
            \item  Si $f : G \to H$ es un homomorfismo de grupos y $N \trianglelefteq G$ es un subgrupo normal, entonces $f_*(N) \trianglelefteq H$ es un subgrupo normal.
            \item  Si $H \trianglelefteq N$ y $N \trianglelefteq G$ son subgrupos normales, entonces $H \trianglelefteq G$ es un subgrupo normal.
            \item  Si $H$ es un subgrupo normal de $G$ y $x^2 \in H$, para todo $x \in G$ entonces $G/H$ es abeliano.
            \item  De $D_{12}$ en $\bb{Z}_{12}$ existen 12 homomorfismos de grupos sobreyectivos.
            \item  Sean $A$ y $B$ dos grupos y consideremos el grupo $G = A \times B$ y un subgrupo normal suyo $H$, tal que $G = HA = HB$ y $H \cap A = H \cap B = 1$, entonces $A$ es isomorfo a $B$.
            \item  Dado un grupo $G = KH$, donde $K, H$ son subgrupos abelianos de $G$ y $K$ es normal en $G$, entonces $G$ es resoluble.
            \item  Si $G$ es un grupo cíclico de orden $p^n$ con $p$ primo, entonces la longitud es $n$ y sus factores son todos $\bb{Z}_p$.
            \item  Si $G$ un grupo abeliano cíclico. Entonces $G$ es simple.
            \item  El centralizador $C_{D_4}(s)$ es isomorfo a $C_2 \times C_2$.
            \item  $A_4 \times \bb{Z}_5$ tiene un único 3-subgrupo de Sylow.
        \end{enumerate}
    \end{ejercicio}

    \begin{ejercicio}[0.5 puntos]
        Sea $G$ el grupo con 24 elementos descrito de la siguiente manera:
        \[
        G = \langle a, b ; a^8 = 1, b^3 = 1, bab = a \rangle
        \]
        \begin{enumerate}[label=\alph*)]
            \item (0.10 puntos) Calcula el orden de $ab$.
            \item (0.10 puntos) ¿Es el subgrupo $H = \langle ab \rangle$ normal?
            \item (0.15 puntos) ¿Tiene $G$ un 2-subgrupo de Sylow normal?
            \item (0.15 puntos) Considera el $G$-conjunto $X = \text{Syl}_2(G)$ y sea $\phi$ el homomorfismo asociado a la acción. Demuestra que la acción no es fiel y que $\ker(\phi)$ no puede tener orden primo.
        \end{enumerate}
    \end{ejercicio}

    \begin{ejercicio}[0.5 puntos]
        Sea $G = D_7 \times H$ donde $H$ es un grupo simple de orden 168. ¿Cuántos elementos de orden 7 tiene $G$?
    \end{ejercicio}

    \newpage
    \setcounter{ejercicio}{0} % Reiniciar contador de ejercicios
    \noindent
    \textbf{Solución.}

    \begin{ejercicio}[1.0 punto]
        \textbf{Cuestionario}
        \begin{enumerate}[label=\arabic*)]
            \item \textbf{FALSO.} 
            Consideramos la cadena $D_2 \trianglelefteq V_4 \trianglelefteq S_4$, 
            donde\\
            $V_4 = \{id, (1\,2)(3\,4), (1\,3)(2\,4), (1\,4)(2\,3)\}$ es el subgrupo normal de Klein en $S_4$ 
            y tomamos $N = \{id, (1\,2)(3\,4)\} \cong D_2$ aunque valdría cualquier generador.

            Al ser $V_4$ abeliano, todo subgrupo suyo es normal, por lo que $N \trianglelefteq V_4$. 
            Si definimos el homomorfismo de inclusión $i: V_4 \hookrightarrow S_4$, la imagen directa de $N$ es $i_*(N) = N$. 
            
            Sin embargo, $N$ no es normal en el codominio $H = S_4$, puesto que al conjugar por la transposición $(2\,3) \in S_4$ 
            obtenemos $(2\,3)(1\,2)(3\,4)(2\,3)^{-1} = (1\,3)(2\,4) \notin N$. 
            Por tanto, la imagen de un subgrupo normal mediante un homomorfismo no es necesariamente normal en el grupo de destino.
            
            \item \textbf{FALSO.}
            Si utilizamos el mismo contraejemplo de antes $D_2 \trianglelefteq V_4 \trianglelefteq S_4$ ya hemos probado que la normalidad no es transitiva.
            
            \item \textbf{VERDADERO.} 
            En el grupo cociente $G/H$, tomemos dos clases cualesquiera $xH$ y $yH$. 
            Por hipótesis, para todo $g \in G$ se cumple que $g^2 \in H$, lo que significa que en el cociente $(gH)^2 = g^2H = H$. 
            De este modo, cada elemento coincide con su propio elemento inverso ($(xH)^{-1} = xH$). Todo grupo en el cual el cuadrado de cada elemento es la identidad es abeliano, ya que:
            \[ xHyH = (xHyH)^{-1} = (yH)^{-1}(xH)^{-1} = yHxH \]
            Por lo tanto, $G/H$ es abeliano.
            
            \item \textbf{FALSO.}
            El número de homomorfismos sobreyectivos es 0.

            Sea $\phi: D_{12} \to \mathbb{Z}_{12}$ un homomorfismo de grupos. Supongamos, por reducción al absurdo, que $\phi$ es sobreyectivo, es decir, $\text{Im}(\phi) = \mathbb{Z}_{12}$.

            Por el Primer Teorema de Isomorfía para grupos, sabemos que:
            \[
            \frac{D_{12}}{\ker(\phi)} \cong \text{Im}(\phi) = \mathbb{Z}_{12}
            \]
            Dado que el grupo cociente $D_{12}/\ker(\phi)$ es isomorfo a $\mathbb{Z}_{12}$, en particular es un grupo abeliano. Por las propiedades de los grupos cocientes abelianos, el núcleo $\ker(\phi)$ debe contener necesariamente al subgrupo conmutador (o derivado) de $D_{12}$, denotado por $D_{12}' = [D_{12}, D_{12}]$.

            Recordemos que la presentación canónica del grupo diédrico $D_{12}$ (de orden $2 \cdot 12 = 24$) es:
            \[
            D_{12} = \langle r, s \mid r^{12} = 1, s^2 = 1, srs = r^{-1} \rangle
            \]
            Calculamos el conmutador genérico de los generadores:
            \[
            [r, s] = r^{-1}s^{-1}rs = r^{-1}(srs) = r^{-1}r^{-1} = r^{-2}
            \]
            Por tanto, el subgrupo derivado está generado por las potencias pares de la rotación: $D_{12}' = \langle r^2 \rangle$. Como $r$ tiene orden 12, el elemento $r^2$ tiene orden $\frac{12}{\text{mcd}(2,12)} = 6$. Así, $|D_{12}'| = 6$.

            El cociente por el subgrupo derivado (la abelización de $D_{12}$) tiene orden $\frac{24}{6} = 4$

            Por el Tercer Teorema de Isomorfía, como $D_{12}' \le \ker(\phi)$, el grupo cociente $D_{12}/\ker(\phi)$ es isomorfo a un subgrupo de $D_{12}/D_{12}'$. 

            En ese caso:
            \[
            |\text{Im}(\phi)| = \left| \frac{D_{12}}{\ker(\phi)} \right| \text{ divide a } |\mathbb{Z}_2 \times \mathbb{Z}_2| = 4
            \]
            En particular, $|\text{Im}(\phi)| \le 4$. Esto contradice la hipótesis de que $\phi$ es sobreyectivo, lo cual requeriría que $|\text{Im}(\phi)| = |\mathbb{Z}_{12}| = 12$. 

            Concluimos que no existe ningún homomorfismo sobreyectivo de $D_{12}$ en $\mathbb{Z}_{12}$.

            \item \textbf{VERDADERO.} 
            Haciendo uso del Segundo Teorema de Isomorfía y apoyándonos en que $H \trianglelefteq G$, podemos relacionar los subgrupos del producto directo de la siguiente manera:
            \[ G/H = (HA)/H \cong A/(H \cap A) = A/1 \cong A \]
            Efectuando el mismo procedimiento de forma simétrica para el subgrupo $B$:
            \[ G/H = (HB)/H \cong B/(H \cap B) = B/1 \cong B \]
            Al ser tanto $A$ como $B$ isomorfos al mismo grupo cociente $G/H$, por la propiedad transitiva de los isomorfismos de grupos se deduce directamente que $A \cong B$.
            
            \item \textbf{VERDADERO.} 
            Consideremos la serie de subgrupos $1 \trianglelefteq K \trianglelefteq G$.
            Analicemos sus factores: el primer factor es $K/1 \cong K$, que es abeliano por hipótesis. 
            El segundo factor es $G/K = (KH)/K$, el cual, aplicando el Segundo Teorema de Isomorfía, resulta ser isomorfo a $H/(K \cap H)$. Dado que $H$ es abeliano, cualquier cociente suyo también conserva la propiedad conmutativa. 
            Al haber construido una serie normal cuyos factores cocientes son todos abelianos vemos que $G$ es un grupo resoluble.
            
            \item \textbf{VERDADERO.} 
            Un grupo cíclico $G \cong \bb{Z}_{p^n}$ posee un único subgrupo de orden $p^k$ para cada entero $0 \le k \le n$.
            Debido a esta unicidad, cuenta con una sola serie de composición, la cual viene determinada por la cadena ordenada de subgrupos $1 = G_0 \trianglelefteq G_1 \trianglelefteq \dots \trianglelefteq G_n = G$, con $|G_k| = p^k$.
            La longitud de la serie coincide con el exponente $n$. Asimismo, cada factor de la composición $G_k/G_{k-1}$ tiene un orden igual a $p^k/p^{k-1} = p$; al ser de orden primo, todos los factores son necesariamente isomorfos a $\bb{Z}_p$.

            También podríamos haber aplicado que al ser cíclico es abeliano y usando el ejercicio \textbf{2.7.5} de la relación de ejercicios sacamos inmediatamente lo pedido.
            
            \item \textbf{FALSO.} 
            Sería cierto si además fuera de orden primo, por tanto, si consideramos el grupo $\bb{Z}_4$ vemos que es abeliano y cíclico pero $\bb{Z}_2 \triangleleft \bb{Z}_4$ por lo que $\bb{Z}_4$ no es simple.
            
            \item \textbf{VERDADERO.} El grupo diédrico $D_4$ se define mediante la presentación $\langle r, s \mid r^4=1, s^2=1, srs=r^{-1}\rangle$. El centralizador $C_{D_4}(s)$ está formado por todos los elementos de $D_4$ que conmutan con la reflexión $s$. De entrada, es evidente que $1, s \in C_{D_4}(s)$. Por otra parte, sabemos que el centro del grupo es $Z(D_4) = \{1, r^2\}$, lo que asegura que el elemento $r^2$ también conmuta con $s$. Por último, el elemento $sr^2$ cumple $s(sr^2) = r^2 = (sr^2)s$, de modo que también pertenece al conjunto. Así, el centralizador es $C_{D_4}(s) = \{1, r^2, s, sr^2\}$, que posee orden 4. Puesto que todos sus elementos no triviales tienen orden 2 (ya que $(r^2)^2=1$, $s^2=1$ y $(sr^2)^2 = sr^2sr^2 = ssr^{-2}r^2 = 1$), este grupo es isomorfo al grupo de Klein, $C_2 \times C_2$.
            
            \item \textbf{FALSO.} El orden del producto directo es:
            $$|A_4 \times \bb{Z}_5| = |A_4| \cdot |\bb{Z}_5| = 12 \cdot 5 = 60 = 2^2 \cdot 3 \cdot 5$$ Los 3-subgrupos de Sylow tienen orden 3. Utilizando los teoremas de Sylow, el número de estos subgrupos, $n_3$, debe dividir a $\nicefrac{60}{3} = 20$ y verificar la condición de congruencia $n_3 \equiv 1 \pmod 3$. Los divisores de 20 son $\{1, 2, 4, 5, 10, 20\}$, y los únicos que cumplen la congruencia son 1 y 4. En el factor $A_4$ existen exactamente 8 elementos de orden 3 (los 3-ciclos), los cuales dan lugar a $\nicefrac{8}{2} = 4$ subgrupos de orden 3. Los 3-subgrupos de Sylow de $A_4 \times \bb{Z}_5$ son de la forma $P \times \{0\}$, con $P \in \text{Syl}_3(A_4)$. Al existir 4 elecciones posibles para $P$, el grupo total cuenta con 4 subgrupos de Sylow de orden 3, de modo que no es único.
        \end{enumerate}
    \end{ejercicio}

    \begin{ejercicio}[0.5 puntos]
        Sea $G$ el grupo con 24 elementos descrito de la siguiente manera:
        \[
        G = \langle a, b ; a^8 = 1, b^3 = 1, bab = a \rangle
        \]
        \begin{enumerate}[label=\alph*)]
            \item Calcula el orden de $ab$.\\
            Para calcular el orden de $ab$, buscaremos el menor entero positivo $k$ tal que $(ab)^k = 1$. Empecemos calculando las potencias de este elemento, aprovechando la relación del enunciado $bab = a$:
            \begin{align*}
                (ab)^2 &= (ab)(ab) = a(bab) = a(a) = a^2 \\
                (ab)^4 &= \left((ab)^2\right)^2 = (a^2)^2 = a^4 \\
                (ab)^8 &= \left((ab)^4\right)^2 = (a^4)^2 = a^8 = 1
            \end{align*}
            Sabemos que $(ab)^8 = 1$, por lo que el orden de $ab$ debe dividir a 8 (es decir, puede ser 1, 2, 4 u 8).
            \begin{itemize}
                \item El orden de $a$ es 8
                \item Por tanto, $a^2 \neq 1$ y $a^4 \neq 1$.
                \item Como $(ab)^2 = a^2 \neq 1$ y $(ab)^4 = a^4 \neq 1$, concluimos que el orden no puede ser ni 2 ni 4.
            \end{itemize}
            \textbf{Conclusión:} El orden de $ab$ es \textbf{8}.
            
            \item ¿Es el subgrupo $H = \langle ab \rangle$ normal?\\
            Sabemos por el apartado anterior que $H = \langle ab \rangle$ es un subgrupo de orden 8. 
            Para comprobar si es normal, basta ver si es invariante por conjugación frente a los generadores de $G$. Conjugaremos el generador de $H$, que es $ab$, mediante el elemento $b$.

            Primero, manipulamos la relación $bab = a$. Multiplicando por $b^{-1}$ por la derecha obtenemos:
            \[
            ba = ab^{-1} = ab^2
            \]
            Ahora conjugamos $ab$ por $b$:
            \[
            b (ab) b^{-1} = (bab)b^{-1} = a b^{-1} = a b^2
            \]
            Para que $H \triangleleft G$, el elemento $ab^2$ debería pertenecer a $H$. Vamos a demostrar que esto es imposible.
            Cualquier elemento de $H$ es de la forma $(ab)^k$. 
            \begin{itemize}
                \item Si $k$ es par, $(ab)^{2m} = a^{2m}$. Si $ab^2 = a^{2m}$, entonces $b^2 = a^{2m-1}$. Esto implicaría que $b^2 \in \langle a \rangle$, lo cual es imposible pues los órdenes de $a$ y $b$ son coprimos ($\text{mcd}(8,3)=1$) y la intersección de sus subgrupos generados es trivial.
                \item Si $k$ es impar, $(ab)^{2m+1} = a^{2m}ab = a^{2m+1}b$. Si $ab^2 = a^{2m+1}b$, cancelando por la izquierda y derecha obtendríamos $b = a^{2m} \in \langle a \rangle$, llegando a la misma contradicción.
            \end{itemize}
            Como $ab^2 \notin H$, tenemos que $b H b^{-1} \neq H$. 
            \\
            \textbf{Conclusión:} El subgrupo $H = \langle ab \rangle$ \textbf{no es normal} en $G$.

            \item  ¿Tiene $G$ un 2-subgrupo de Sylow normal?\\
            El orden del grupo es $|G| = 24 = 2^3 \cdot 3$, lo que nos indica que sus 2-subgrupos de Sylow tienen un orden de $2^3 = 8$. 
            El subgrupo $H = \langle ab \rangle$ estudiado en los apartados anteriores tiene orden 8, clasificándose como un 2-subgrupo de Sylow de $G$. 
            Por los teoremas de Sylow, sabemos que todos los subgrupos de Sylow para un primo dado son conjugados entre sí. 
            Si existiese un 2-subgrupo de Sylow normal, este tendría que ser único ($n_2=1$) y coincidir necesariamente con $H$. Habiendo demostrado en el apartado (b) que $H$ no es normal, vemos claramente que $G$ \textbf{no tiene ningún 2-subgrupo de Sylow normal}.
            
            \item  Considera el $G$-conjunto $X = \text{Syl}_2(G)$ y sea $\phi$ el homomorfismo asociado a la acción. Demuestra que la acción no es fiel y que $\ker(\phi)$ no puede tener orden primo.\\
            Consideremos la acción por conjugación del grupo $G$ sobre el conjunto de sus 2-subgrupos de Sylow, $X = \text{Syl}_2(G)$, donde $|X| = n_2 = 3$. 
            
            Sea $\phi: G \to S_X \cong S_3$ el homomorfismo de permutaciones asociado a dicha acción. El núcleo de este homomorfismo se define como la intersección de todos los 2-subgrupos de Sylow: $\ker(\phi) = \bigcap_{P \in \text{Syl}_2(G)} P$. Pues recordemos que todos son conjugados entre sí.
            
            Por el Primer Teorema de Isomorfía, el grupo cociente $G/\ker(\phi)$ es isomorfo a un subgrupo del grupo simétrico $S_3$. Aplicando el Teorema de Lagrange, el orden del cociente debe dividir al orden de $S_3$, es decir, a $6$:
            \[ |G/\ker(\phi)| = \frac{24}{|\ker(\phi)|} \ \Big| \ 6 \implies |\ker(\phi)| \ge \frac{24}{6} = 4 \]
            Dado que el orden del núcleo es $|\ker(\phi)| \ge 4 > 1$, este contiene elementos distintos de la identidad, lo que prueba rigurosamente que la acción \textbf{no es fiel}.
            
            Por otra parte, aplicando de nuevo el Teorema de Lagrange, el orden del subgrupo $\ker(\phi)$ tiene que ser un divisor del orden total del grupo, $|G|=24$. Los divisores de 24 que cumplen con la condición de ser mayores o iguales que 4 son los elementos del conjunto $\{4, 6, 8, 12, 24\}$. Como se puede observar, ninguno de estos valores corresponde a un número primo, demostrando así que el orden de $\ker(\phi)$ \textbf{no puede ser primo}.
        \end{enumerate}
    \end{ejercicio}

    \begin{ejercicio}[0.5 puntos] Sea $G = D_7 \times H$ donde $H$ es un grupo simple de orden 168. ¿Cuántos elementos de orden 7 tiene $G$?\\
        Dado cualquier $(x,y) \in D_7 \times H$ sabemos que $O(x, y) = \mcm(O(x), O(y))$. Debido a que 7 es un número primo, esto pasa si y solo si se da uno de los siguientes tres casos:
        \begin{itemize}
            \item \textbf{Caso 1:} $|x| = 7$ y $|y| = 1$.
            \item \textbf{Caso 2:} $|x| = 1$ y $|y| = 7$.
            \item \textbf{Caso 3:} $|x| = 7$ y $|y| = 7$.
        \end{itemize}
        
        Evaluemos de forma independiente la cantidad de elementos de orden 7 en cada factor, pues de orden 1 solo está la identidad:
        \begin{itemize}
            \item \textbf{En el factor $D_7$:} Sabemos que $C_7\cong \langle r\rangle \le D_7$ y como es cíclico de orden primo, todos sus elementos son de orden 7. El resto de elementos son las simetrías que sabemos que todas tienen orden 2.
            Por tanto en total tenemos 6 elementos de orden 7 y la identidad con orden 1. 
            
            \item \textbf{En el factor $H$:} El grupo $H$ es simple y tiene un orden de $168 = 2^3 \cdot 3 \cdot 7$.
            Los elementos de orden 7 se encuentran contenidos en los 7-subgrupos de Sylow de $H$.
            Como dichos subgrupos tienen orden primo 7, su intersección es trivial y cada uno de ellos aporta exactamente $7 - 1 = 6$ elementos de orden 7.
            
            Determinamos el número $n_7$ de 7-subgrupos de Sylow mediante las condiciones de Sylow:
            \begin{align*}
                n_7 \equiv 1 \pmod 7\\
                n_7 \mid \frac{168}{7} = 24
            \end{align*}
            Los divisores de 24 son $\{1, 2, 3, 4, 6, 8, 12, 24\}$, de los cuales solo el 1 y el 8 satisfacen la condición de congruencia. 
            Si $n_7 = 1$, el único 7-subgrupo de Sylow sería normal en $H$, lo que contradice directamente la hipótesis de que $H$ es un grupo simple. 
            En consecuencia, el número de subgrupos de Sylow es obligatoriamente $n_7 = 8$. Así, el número total de elementos de orden 7 en $H$ es $8 \cdot 6 = 48$ elementos y la identidad con orden 1.
        \end{itemize}
        
        Calculamos el número de combinaciones posibles para cada uno de los escenarios planteados:
        \begin{itemize}
            \item \textbf{Caso 1:} $6 \text{ (de } D_7) \cdot 1 \text{ (de } H) = 6$ elementos.
            \item \textbf{Caso 2:} $1 \text{ (de } D_7) \cdot 48 \text{ (de } H) = 48$ elementos.
            \item \textbf{Caso 3:} $6 \text{ (de } D_7) \cdot 48 \text{ (de } H) = 288$ elementos.
        \end{itemize}
        Efectuando la suma de las posibilidades de todos los casos concurrentes, obtenemos que el número total de elementos de orden 7 en el grupo $G$ es:
        \[ 6 + 48 + 288 = 342 \]
    \end{ejercicio}

\end{document}